\section{Study Design}
\label{sec:study-design}

\subsection{Research Questions}
\label{sec:rq}

The research questions were frozen before outcome generation and are reproduced below:

\begin{description}
  \item[RQ1.] How stable are scheduler rankings across independent
    workload sources?
  \item[RQ2.] How stable are scheduler rankings under temporal, provider,
    and domain shifts?
  \item[RQ3.] To what extent do rankings obtained on synthetic stress
    workloads transfer to rankings on independent real-trace-derived
    workloads?
  \item[RQ4.] Which source-native observable workload characteristics are
    associated with cross-distribution scheduler rank reversals?
  \item[RQ5.] How many independent workload windows are required before a
    comparative scheduler ranking becomes statistically stable?
  \item[RQ6.] Do representative simulated scheduler rankings and
    cross-workload rank reversals reproduce on a real serving engine?
\end{description}

Load-level dependence is a secondary robustness/sensitivity analysis applied
to RQ1--RQ4's findings, not a headline research question
(\S\ref{sec:evidence-boundaries}).

\subsection{Study Overview}
\label{sec:overview}

LSSP evaluates a fixed panel of 13 scheduling policies (\S\ref{sec:panel})
against a frozen corpus of 120 workload windows drawn from three
independently sourced traces (\S\ref{sec:sources}), across a calibrated
six-region grid of operating regions (\S\ref{sec:calibration}) and a fixed
set of primary and completion-conditioned metrics (\S\ref{sec:metrics}).
Every policy is run against every window under every operating region,
with a deterministic rep0/rep1 verification pass (\S\ref{sec:paired}),
producing $9{,}360$ unique policy--window--region configurations, each executed in two deterministic verification passes ($18{,}720$ total executions), from which per-source, per-region, and
per-metric scheduler rankings are derived. Ranking-portability statistics (Kendall's tau-b, Spearman's
rho, top-\emph{k} overlap, and pairwise reversal detection) are then
computed across pairs of sources, regions, and metrics, using workload
window as the statistical unit (\S\ref{sec:unit}). This design was frozen
(\S\ref{sec:prereg}) before any comparative scheduler outcome existed.

\subsection{Evidence Independence / Publication Boundaries}
\label{sec:evidence-boundaries}

This project shares a lineage and a simulator codebase with a related
manuscript that could otherwise appear to overlap in evidence, and its
design was explicitly checked against it for overlap before the protocol
was frozen:

\begin{itemize}
  \item \textbf{``The Exploitability Gap in LLM-Serving Scheduler
    Portfolios'' (LLM 2026).} That manuscript's own experiment used a
    fixed 60-window subset of public BurstGPT and Azure-2023 traces, a
    different load-factor grid, and an 8-policy portfolio; it is cited
    here strictly as prior motivation and never restated as this
    project's own evidence. This project's window sampling
    (\S\ref{sec:windows}) was independently constructed and verified
    non-identical to that 60-window set, and its load-level dependence
    analysis (\S\ref{sec:rq}) uses this project's own calibration
    protocol (\S\ref{sec:calibration}), never that load grid.
  \item \textbf{A separate module-intervention manuscript (``Module
    Counterfactuals for Attributing LLM Serving Scheduler Effects'').}
    This study makes no module-counterfactual or module-attribution
    claim; RQ4's workload-descriptor analysis (\S\ref{sec:methodology})
    is an offline, descriptive association analysis, not a deployed
    selector or an intervention-based attribution method.
\end{itemize}

\subsection{Experimental Unit and Paired Design}
\label{sec:unit}

The statistical unit of analysis is the workload window: a fixed-size
(200-request) contiguous slice of a source trace, uniquely identified and
frozen before any scheduler outcome existed (\S\ref{sec:windows}). Every
policy in the panel (\S\ref{sec:panel}) is run against every window under
every calibrated operating region (\S\ref{sec:calibration}), so that
ranking-portability statistics are computed on genuinely paired
observations -- the same window, region, and metric definition, varying
only the policy -- rather than on separately tuned per-source experiments.
Source, operating region, and metric are treated as factors of interest;
policy identity is the object whose relative ordering is measured.

\subsection{Preregistration}
\label{sec:prereg}

The workload window set (120 windows: 40 per source, 10 reused from the
earlier discriminability pilot plus 30 newly sampled per source;
\S\ref{sec:windows}), the scheduler panel and
its selection rule (\S\ref{sec:panel}), the metric contract
(\S\ref{sec:metrics}), the six-region operating-load calibration
(\S\ref{sec:calibration}), and the ranking-portability analysis plan
(reversal-detection thresholds, bootstrap procedure, sample-complexity
design) were each frozen, and structurally validated, before any
comparative scheduler outcome was generated. The exact frozen content is
recorded in the public artifact's protocol documentation
(\S\ref{sec:artifact}).
